\section{\texorpdfstring{$\Z/2$}{Z/2} 作用}\label{sec:6}
\subsection{基本概念}\label{sec:6a}
本节讨论带有最简单意义下的对合的 $\Ai$ 范畴。
\begin{definition*}
$\Ai$ 范畴 $\A$ 上的\emph{朴素 $\Z/2$ 作用}是一个 $\Ai$ 函子 $\epsilon_{\A}:\A\to\A$，其阶数不小于二的项为零，并满足 $\epsilon_{\A}\circ\epsilon_{\A}=\Id_{\A}$。
\end{definition*}
事实上，它只由对象上的对合与相应的态射映射 $\epsilon_{\A}:\hom_{\A}(X_0,X_1)\to\hom_{\A}(\epsilon_{\A}X_0,\epsilon_{\A}X_1)$ 组成，满足
\[
\mu^d_{\A}(\epsilon_{\A}(a_d),\ldots,\epsilon_{\A}(a_1))=\epsilon_{\A}(\mu^d_{\A}(a_d,\ldots,a_1)).
\]
\begin{definition*}
若 $\epsilon_{\A}X=X$，就称 $X$ 为 $\epsilon_{\A}$-\emph{不动对象}。\emph{不动 $\Ai$ 子范畴} $\A^{\mathrm{fix}}$ 是由这些对象组成的全 $\Ai$ 子范畴。\emph{不变 $\Ai$ 子范畴} $\A^{\mathrm{inv}}$ 具有与 $\A^{\mathrm{fix}}$ 相同的对象，但其态射为不变部分 $\hom_{\A^{\mathrm{inv}}}(X_0,X_1)=\hom_{\A}(X_0,X_1)^{\Z/2}$，复合映射为相应的限制。
\end{definition*}
我们需要取不变部分这一过程具有良好的同调性质，因此从现在起假设 $\operatorname{char}(\K)\ne2$。

后面会多次遇到以下初等问题。令 $A=H^0(\A)$，$\epsilon_A$ 为 $\epsilon_{\A}$ 所诱导的 $\Z/2$ 作用。取两个 $\epsilon_{\A}$-不动对象 $X_0,X_1$，假设它们在 $A$ 中同构。何时可以把这个同构
\SourcePage{96}
选为 $\Z/2$-不变的？障碍是类
\begin{equation}\label{eq:6.1}
[y=\epsilon_A(x^{-1})\cdot x]\in\{y\in\operatorname{Aut}_A(X_0):y^{-1}=\epsilon_A(y)\}/\operatorname{Aut}_A(X_0),
\end{equation}
其中 $x$ 为任意一个从 $X_0$ 到 $X_1$ 的同构，而 $\operatorname{Aut}_A(X_0)=\Hom_A^0(X_0,X_0)^{\times}$ 通过扭曲共轭作用于自身，即 $z$ 将 $y$ 映为 $\epsilon_A(z)^{-1}yz$。\znote{88}{这个障碍类可以化归为纯代数验算：若 $x:X_0\to X_1$ 为同构，令 $y=\epsilon_A(x^{-1})x$，则 $y^{-1}=x^{-1}\epsilon_A(x)=\epsilon_A(y)$，故 $y$ 必落在所示集合中；若换用另一同构 $xz$（$z\in\operatorname{Aut}_A(X_0)$），$y$ 就变成 $\epsilon_A(z)^{-1}yz$。而 $y$ 为单位元恰等价于 $\epsilon_A(x)=x$。因此这个商集退化为单点，当且仅当存在 $\Z/2$-不变的同构；扭曲共轭正是把「换一个 $x$」造成的差别除掉的等价关系。[来源：译者按（上下文推断，无外部依据）：由 $y$ 的定义直接验算 $y^{-1}=\epsilon_A(y)$，并代入 $x\mapsto xz$ 得到扭曲共轭公式。]}
\begin{definition*}
最简单的情形是 $X_0$ \emph{自由}：这意味着 $\Hom_A(X_0,X_0)\cong R\oplus R$ 是同一个分次代数 $R$ 的两个副本的直和，且 $\epsilon_A$ 通过交换两个直和项作用。
\end{definition*}
在这种情形下，\eref{6.1}中的集合由 $y=(r,r^{-1})$ 这样的元素组成，每个元素都可以通过以 $z=(r^{-1},e)$ 作扭曲共轭而化为平凡元素。因此：
\begin{lemma}\label{st:6.1}
设有两个 $\epsilon_{\A}$-不动对象，其中一个已知是自由的。那么，若它们在 $H^0(\A)$ 中同构，它们也在 $H^0(\A^{\mathrm{inv}})$ 中同构。
\end{lemma}

\subsection{等变扭曲复形}\label{sec:6b}
首先需要讨论加性扩张的等变版本 $\Sigma^{\mathrm{eq}}\A$。其对象是 $\Sigma\A$ 的对象，即写成\eref{3.15}中形式和的三元组 $(I,\{X^i\},\{V^i\})$，再加上下列数据：指标集 $I$ 上的对合 $i\mapsto\underline i$，满足 $X^{\underline i}=\epsilon_{\A}(X^i)$；以及向量空间的分次映射 $\epsilon_{V^i}:V^i\to V^{\underline i}$，对所有 $i$ 满足 $\epsilon_{V^{\underline i}}\circ\epsilon_{V^i}=\Id$。态射空间与复合映射的定义与非等变版本相同。每个态射空间带有诱导的对合；用\eref{3.16}的记号，它由以下两个映射的张量积得到：
\[
\epsilon_{V_1^j}\circ-\circ\epsilon_{V_0^{\underline i}}:
\Hom_{\K}(V_0^i,V_1^j)\longrightarrow\Hom_{\K}(V_0^{\underline i},V_1^{\underline j})
\]
以及
\[
\epsilon_{\A}:\hom_{\A}(X_0^i,X_1^j)\longrightarrow\hom_{\A}(X_0^{\underline i},X_1^{\underline j}).
\]
\begin{definition*}
\emph{等变扭曲复形}是二元组 $(X,\delta_X)$，其中 $X\in\Ob\Sigma^{\mathrm{eq}}\A$，微分 $\delta_X$ 属于不变部分 $\hom^1_{\Sigma^{\mathrm{eq}}\A}(X,X)^{\Z/2}$，并满足通常的条件：存在滤过（此时应由等变子复形 $F^iX$ 组成），以及 Maurer–Cartan 方程。
\end{definition*}
态射空间沿用 $\Sigma^{\mathrm{eq}}\A$ 中的定义，复合映射按\eref{3.20}定义，所得为 $\Ai$ 范畴 $\Tw^{\mathrm{eq}}\A$，它带有在对象上平凡的朴素 $\Z/2$ 作用。

最简单的等变扭曲复形是 $\A$ 中的单个 $\epsilon_{\A}$-不动对象 $X$（更精确地说，取固定的单元素集 $I=\{*\}$、一维向量空间 $V^*=\K$，其映射 $\epsilon_{V^*}=\Id$ 为平凡映射，并取 $X^*=X$）。一个稍不平凡的例子是 $\A$ 中任意对象与其像的直和 $X\oplus\epsilon_{\A}X$；微分仍为零，但对合 $i\mapsto\underline i$ 交换两个直和项。注意，第一类对象给出 $\A^{\mathrm{fix}}$ 到 $\Tw^{\mathrm{eq}}\A$ 的全嵌入，与 $\Z/2$ 作用相容。更一般地，$\A^{\mathrm{inv}}$ 中的任意扭曲复形都可以
\SourcePage{97}
视为等变扭曲复形（集 $I$ 与向量空间 $V^i$ 上的对合均平凡）；这样，$\Tw(\A^{\mathrm{inv}})$ 成为 $(\Tw^{\mathrm{eq}}\A)^{\mathrm{inv}}$ 的全 $\Ai$ 子范畴。最后，直接忘掉对合，便定义了全忠实 $\Ai$ 函子 $\Tw^{\mathrm{eq}}\A\to\Tw\A$。

再说明一些熟悉运算的等变版本。直和的定义是显然的。其次，给定带 $\Z/2$ 作用 $\epsilon_Z$ 的有限维分次向量空间 $Z$ 和等变扭曲复形 $X$，可以取张量积 $Z\otimes V^i$，赋予组合的 $\Z/2$ 作用 $\epsilon_Z\otimes\epsilon_{V^i}$，并取微分 $\delta_Y=\id_Z\otimes\delta_X$，从而形成等变张量积 $Y=Z\otimes X$。这包括平移运算 $S$，其定义照常是取带平凡 $\Z/2$ 作用的 $Z=\K[1]$。另一个有趣情形是 $\Z/2$ 的非平凡特征标，即 $Z=\K$、$\epsilon_Z=-\id_Z$。这时张量积只是将所有 $\epsilon_{V^i}$ 的符号反转，把这一过程记为 $X\mapsto\check X$。一般的张量积 $Z\otimes X$ 本质上是 $X$ 与 $\check X$ 的若干平移的有限和，$Z$ 的每个不变或反不变基元素各对应一项。最后，若 $c\in\hom_{\Tw\A}(Y_0,Y_1)$ 是不变的零次上闭链，那么赋予 $SY_0\oplus Y_1$ 微分\eref{3.28}所定义的映射锥，仍是等变扭曲复形。这说明 $(\Tw^{\mathrm{eq}}\A)^{\mathrm{inv}}$ 三角化，而 $\Tw^{\mathrm{eq}}\A$ 本身并不三角化。\znote{89}{成败取决于特征标：等变扭曲复形的每个对象都带有 $\Z/2$ 在其上的作用，可以平凡，也可以按 $-\id$ 作用（后者即 $\check X$）。要在 $SY_0\oplus Y_1$ 上写出对合，必须让两块的符号一致，这只有当 $c$ 落在不变部分时才做得到；对一般的 $c$，锥没有自然的等变结构。反过来，一旦取不变部分，$Y_0,Y_1,c$ 全都不变，锥便自动留在 $(\Tw^{\mathrm{eq}}\A)^{\mathrm{inv}}$ 里——这就是「只有取不变部分之后才三角化」的确切含义。[来源：原书 §6b 本段关于锥的构造与 (3.28) 的锥定义，并见本段关于 $\check X$ 的讨论。]}

\begin{lemma}\label{st:6.2}
设 $Y,Y_0,Y_1\in\Ob\A$ 均为 $\epsilon_{\A}$-不动对象；另假设 $Y$ 自由，且 $\Hom_{H(\A)}(Y_0,Y_1)$ 上的 $\Z/2$ 作用平凡。若在 $D(\A)=H^0(\Tw\A)$ 中，$Y$ 同构于某个态射 $Y_0\to Y_1$ 的锥，那么同样的结论在 $D(\A^{\mathrm{inv}})$ 中成立。
\end{lemma}
\begin{proof}
由假设，$\Hom_{H(\A)}(Y_0,Y_1)$ 是 $\Z/2$-不变的，因此可以在不改变锥的同构类的情况下，假设有关上闭链 $c\in\hom^0_{\A}(Y_0,Y_1)$ 本身 $\Z/2$-不变。这时可以把 $\Cone(c)$ 构造为 $\Tw^{\mathrm{eq}}\A$ 的对象。事实上，它位于嵌入 $\Tw(\A^{\mathrm{inv}})\subset\Tw^{\mathrm{eq}}\A$ 的像中。由于遗忘函子 $H^0(\Tw^{\mathrm{eq}}\A)\to H^0(\Tw\A)$ 全忠实，就在 $H^0(\Tw^{\mathrm{eq}}\A)$ 中得到同构 $Y\cong\Cone(c)$。再由于 $Y$ 自由，利用\stref{引理}{6.1}可以使这个同构不变；这就意味着它来自 $H^0(\Tw(\A^{\mathrm{inv}}))$ 中的同构。
\end{proof}
\begin{lemma}\label{st:6.3}
设 $\A$ c-有限，$X$ 为 $\epsilon_{\A}$-不动对象。假设 $X$ 拟表示某个对象 $Y$ 与其 $\epsilon_{\A}$ 像的直和 $Y\oplus\epsilon_{\A}Y$，并且 $\Hom_{H(\A)}(Y,\epsilon_{\A}Y)=0$。再假设在 $D(\A)$ 中有同构
\begin{equation}\label{eq:6.2}
T_{Y_1}\cdots T_{Y_m}(Y)\cong S(\epsilon_{\A}Y).
\end{equation}
这里 $Y_1,\ldots,Y_m$ 是 $\A$ 的 $\epsilon_{\A}$-不动对象，并且对所有 $j<k$，$\Hom_{H(\A)}(Y_j,Y_k)$ 上的 $\Z/2$ 作用平凡；$S$ 为平移。那么 $X$ 属于 $\Tw(\A^{\mathrm{inv}})$ 中由 $Y_k$ 生成的三角化子范畴。
\end{lemma}
\SourcePage{98}
\begin{proof}
最清楚的情形是 $\A$ 严格含幺且 $\mu^1_{\A}=0$（这尤其意味着 $\hom_{\A}$ 空间本身有限维）。迭代\eref{5.16}，可知\eref{6.2}左边的对象具有形式
\begin{equation}\label{eq:6.3}
\begin{aligned}
C=Y\oplus\bigoplus\bigl(&\hom_{\A}(Y_{j_r},Y)[1]\otimes\hom_{\A}(Y_{j_{r-1}},Y_{j_r})[1]\otimes\cdots\\[-2pt]
&\otimes\hom_{\A}(Y_{j_1},Y_{j_2})[1]\otimes Y_{j_1}\bigr),
\end{aligned}
\end{equation}
其中对所有 $r\geq1$ 与 $j_1<j_2<\cdots<j_r$ 取直和。无须显式写出完整微分 $\delta_C$，只需了解它的一般形式，这很容易由构造推出。首先：
\begin{itemize}
\item $\delta_C$ 严格减小张量积表达式的长度。特别地，它在第一个直和项 $Y$ 上为零。
\end{itemize}
因此，$C$ 可以写为态射 $U\to Y$ 的锥，其中 $(U,\delta_U)\in\Ob\Tw\A$ 包含\eref{6.3}中除首项之外的所有项，整体反向平移一次，并带有微分的相应部分。其次：
\begin{itemize}
\item 只考虑满足 $j_1\geq j$ 的直和项所给的滤过与 $\delta_U$ 相容。此外，相对于这个滤过的对角项，即相伴分次空间上的诱导微分，形如 $\partial_{Z_j}\otimes e_{Y_j}$，其中 $\partial_{Z_j}$ 是重数空间
\[
Z_j=\bigoplus_{\substack{r\geq1\\j=j_1<j_2<\cdots<j_r}}
\bigl(\hom_{\A}(Y_{j_r},Y)\otimes\hom_{\A}(Y_{j_{r-1}},Y_{j_r})\otimes\cdots\otimes\hom_{\A}(Y_j,Y_{j_2})\bigr)[r-1]
\]
上的某个微分；$r=1$ 时该项理解为 $\hom_{\A}(Y_j,Y)$。仔细考察会发现 $\partial_{Z_j}$ 是某种条形微分，但我们不用这一事实。
\end{itemize}
结合 $j<k$ 时 $\hom_{\A}(Y_j,Y_k)$ 完全 $\Z/2$-不变的假设，可知 $\delta_U$ 的每个系数，都是某个向量空间同态与 $\A$ 中某个 $\epsilon_{\A}$-不变态射的张量积。因此，$U$ 属于子范畴 $\Tw(\A^{\mathrm{inv}})\subset\Tw\A$；再看上述描述，可知它实际上属于 $\Tw(\A^{\mathrm{inv}})$ 中由 $Y_k$ 生成的三角化子范畴。根据\eref{6.2}，在 $H(\Tw\A)$ 中有形如
\SourceDiagram{diagram-6-3}{0.24299974}
的正合三角。由假设，在 $H(\A)$ 中，从 $Y$ 到 $\epsilon_{\A}Y$ 没有任何次数的非零态射，因此该正合三角分裂，即在 $H^0(\Tw\A)$ 中有
\begin{equation}\label{eq:6.4}
X\cong Y\oplus\epsilon_{\A}Y\cong U.
\end{equation}
\eref{6.4}中的所有对象都是等变扭曲复形，且中间的对象自由。用\stref{引理}{6.2}证明结尾的论证，便得到 $H^0(\Tw(\A^{\mathrm{inv}}))$ 中所需的同构。

\SourcePage{99}
为将结论推广到一般 $\A$，去掉上面引入的技术限制，作如下处理。选择\stref{笔记}{1.13}中的数据，但要保证复形 $\hom_{\A}(X_0,X_1)$ 上的分裂与同伦和 $\Z/2$ 作用相容。只有 $X_0,X_1$ 都为 $\epsilon_{\A}$-不动对象时，才可能有困难；在这些情形下，只需将复形分成不变与反不变直和项，分别处理后再合起来。于是，\stref{扰动引理}{1.12}之后的公式\eref{1.18}保持原有的对称性，即得到等变拟同构 $\F:\A\to\B$，其中 $\B$ 带有朴素 $\Z/2$ 作用且 $\mu^1_{\B}=0$。接着，像\stref{引理}{2.1}那样使 $\B$ 严格含幺。该引理证明的第一步是平凡的：由于 $\mu^1_{\B}=0$，上同调单位元具有唯一的代表，且自动 $\Z/2$-不变，可以取 $\Phi^2=0$。构造的其余部分不涉及选择，见\eref{2.4}，因而自然与对合相容，给出等变形式微分同胚 $\Phi$，使 $\widetilde{\B}=\Phi_*\B$ 严格含幺。取 $\G=\Phi\circ\F$，考虑 $X,Y,Y_k$ 在其下的像。显然，在 $H^0(\widetilde{\B})$ 中有
\[
\begin{gathered}
\G X\cong\G Y\oplus\G(\epsilon_{\A}Y)\cong\G Y\oplus\epsilon_{\widetilde{\B}}(\G Y),\\
\Hom^*_{H(\widetilde{\B})}(\G Y,\epsilon_{\widetilde{\B}}(\G Y))=0,\\
\Hom^*_{H(\widetilde{\B})}(\G Y_i,\G Y_j)\text{ 在 }i<j\text{ 时为 }\Z/2\text{-不变的}.
\end{gathered}
\]
由\stref{引理}{5.6}，\eref{6.2}的对应式也成立。因此处于上面考虑的特殊情形，即 $\G X$ 属于 $\Tw(\widetilde{\B}^{\mathrm{inv}})$ 中由 $\G Y_k$ 生成的三角化子范畴。但不变部分 $\G^{\mathrm{inv}}:\A^{\mathrm{inv}}\to\widetilde{\B}^{\mathrm{inv}}$ 也是拟同构，因此由\stref{引理}{3.25}诱导导出范畴的等价，所需结论随之成立。
\end{proof}

\subsection{\texorpdfstring{$\Ai$}{A∞} 模}\label{sec:6c}
与上面对扭曲复形的讨论相对应，也可以引入 $\Ai$ 模的等变版本。
\begin{definition*}
$\A$ 上的\emph{等变 $\Ai$ 模}是 $\Ai$ 模 $\M$，连同模同态 $\epsilon_{\M}:\M\to\epsilon_{\A}^*\M$，要求其阶数不小于二的项为零，并且 $\epsilon_{\M}$ 与 $\epsilon_{\A}^*(\epsilon_{\M})$ 以任一顺序复合，均为恒等模自同态。\znote{91}{记号 $\epsilon_{\A}^*\M$ 需要说明：沿对合 $\epsilon_{\A}$ 拉回，把 $\A$-模 $\M$ 变成另一个 $\A$-模，其分量集合仍是 $\M(X)$，但结构映射改为经由 $\M(\epsilon_{\A}X)$ 给出。因此所谓模同态 $\epsilon_{\M}:\M\to\epsilon_{\A}^*\M$，逐分量就是下一句写出的那族映射 $\M(X)\to\M(\epsilon_{\A}X)$，并要求与 $\mu^d_{\M}$ 相容；写成模同态而非逐分量数据，好处是「$\epsilon_{\M}$ 给出等变结构」这句话可以直接用模范畴的语言陈述，不必每次重复相容性公式。[来源：原书 §6c 本段及紧接的分量公式，并见 §6b 中 $\Sigma^{\mathrm{eq}}\A$ 的同类构造。]}
\end{definition*}
具体地，附加数据由一族映射 $\epsilon_{\M}:\M(X)\to\M(\epsilon_{\A}(X))$ 组成，$X$ 与 $\epsilon_{\A}(X)$ 对应的映射互为逆，并且以如下方式与 $\epsilon_{\A}$ 相容：
\[
\mu^d_{\M}(\epsilon_{\M}(b),\epsilon_{\A}(a_{d-1}),\ldots,\epsilon_{\A}(a_1))=\epsilon_{\M}(\mu^d_{\M}(b,a_{d-1},\ldots,a_1)).
\]
态射及其复合按非等变情形定义，将所得 $\Ai$ 范畴记为 $\Q=\modcat^{\mathrm{eq}}(\A)$。它带有诱导的朴素 $\Z/2$
\SourcePage{100}
作用，在对象上平凡；具体地，对 $t\in\hom_{\Q}(\M_0,\M_1)$，令
\[
(\epsilon_{\Q}t)^d(b,a_{d-1},\ldots,a_1)=\epsilon_{\M_1}\bigl(t^d(\epsilon_{\M_0}(b),\epsilon_{\A}(a_{d-1}),\ldots,\epsilon_{\A}(a_1))\bigr).
\]
若 $Y$ 是 $\A$ 的 $\epsilon_{\A}$-不动对象，相应 Yoneda 模 $\Y$ 带有自然的等变结构，即在 $\Y(X)=\hom_{\A}(X,Y)$ 上取 $\epsilon_{\Y}=\epsilon_{\A}$。这定义了上同调全忠实 $\Ai$ 函子 $\mathcal I_{\A}^{\mathrm{eq}}:\A^{\mathrm{fix}}\to\Q$，称为\emph{等变 Yoneda 嵌入}；它与 $\Z/2$ 作用严格交换（利用\stref{引理}{3.36}的等变版本，可以将它延伸为上同调全忠实函子 $\Tw^{\mathrm{eq}}\A\to\Q$）。另一方面，若 $\M$ 为等变 $\Ai$ 模，则只对 $\epsilon_{\A}$-不动对象 $X$ 取空间 $\M(X)$ 的 $\Z/2$-不变部分，就得到 $\A^{\mathrm{inv}}$ 上的 $\Ai$ 模 $\M^{\mathrm{inv}}$。连同对 $\Z/2$-不变自然变换的相应限制，这定义了\emph{不变部分 $\Ai$ 函子} $\Q^{\mathrm{inv}}\to\modcat(\A^{\mathrm{inv}})$。最后，直接忘掉映射 $\epsilon_{\M}$，就得到全忠实 $\Ai$ 函子 $\Q\to\modcat(\A)$。

给定带 $\Z/2$ 作用的向量空间链复形 $Z$ 与等变模 $\M$，张量积 $Z\otimes\M$ 仍然等变。其次，给定两个等变模之间的 $\Z/2$-不变同态，\eref{3.32}意义下的锥仍为等变模。最后，若 $Y$ 是 $\epsilon_{\A}$-不动对象，$\M$ 为等变模，那么扭转模 $\mathcal T_Y\M$ 自然等变，而且在 $H(\Q^{\mathrm{inv}})$ 中有正合三角\eref{5.3}。
\begin{lemma}\label{st:6.4}
设 $Y_0,Y_1,Y_2$ 是 $\epsilon_{\A}$-不动对象，且 $\Hom_{H(\A)}(Y_0,Y_1)$ 上的 $\Z/2$ 作用平凡。假设在 $H^0(\Q)$ 中有
\begin{equation}\label{eq:6.5}
\mathcal T_{Y_0}(\Y_1)\cong\Y_2,
\end{equation}
其中同构属于该范畴态射群的 $\Z/2$-不变部分。那么在 $H^0(\modcat(\A^{\mathrm{inv}}))$ 中有与\eref{6.5}平行的同构，这里的 Yoneda 模与扭转都相对于 $\A^{\mathrm{inv}}$ 而取。最后，若 $\A$ c-有限，则在 $D(\A^{\mathrm{inv}})$ 中有 $T_{Y_0}(Y_1)\cong Y_2$。
\end{lemma}
\begin{proof}
为明确区分各个范畴，将等变 Yoneda $\Ai$ 模记为 $\mathcal I_{\A}^{\mathrm{eq}}(Y)$，将 $\modcat(\A^{\mathrm{inv}})$ 中的对应模记为 $\mathcal I_{\A^{\mathrm{inv}}}(Y)$。应用不变部分函子 $\Q^{\mathrm{inv}}\to\modcat(\A^{\mathrm{inv}})$。它将 $\mathcal I_{\A}^{\mathrm{eq}}(Y_2)$ 映为 $\mathcal I_{\A^{\mathrm{inv}}}(Y_2)$。另一方面，$\mathcal T_{Y_0}(\mathcal I_{\A}^{\mathrm{eq}}(Y_1))$ 的像是 $\A^{\mathrm{inv}}$ 上的 $\Ai$ 模 $\C$，满足
\[
\C(X)=\bigl((\hom_{\A}(Y_0,Y_1)\otimes\hom_{\A}(X,Y_0))[1]\bigr)^{\Z/2}\oplus\hom_{\A}(X,Y_1)^{\Z/2},
\]
并带有通常的结构映射\eref{5.1}。由于 $\hom_{\A}(Y_0,Y_1)$ 的反不变部分无上同调，将第一个直和项替换为 $(\hom_{\A}(Y_0,Y_1)^{\Z/2}\otimes\hom_{\A}(X,Y_0)^{\Z/2})[1]$，就得到一个拟同构子模；由\stref{引理}{1.16}，它在 $H^0(\modcat(\A^{\mathrm{inv}}))$ 中是同构对象。
\SourcePage{101}
但按定义，这个子模恰好就是 $\mathcal T_{Y_0}(\mathcal I_{\A^{\mathrm{inv}}}(Y_1))$。结合\eref{5.6}，最后一项断言立即成立。
\end{proof}
\begin{note}\label{st:6.5}
最后作一些略为宽泛的一般评论。对 $\Ai$ 范畴上群作用采取朴素处理，与同伦代数结构的思想并不一致；但对于有限群以及特征为零、或特征 $p\gg0$ 足够大的基域 $\K$，它仍然相当有效。\secref{10b}会重新考虑一般情形；不过，那里提出的定义虽然比这里的更精细，也至多只是初步近似。\znote{93}{「朴素」在此有确切含义：把 $\Z/2$ 作用实现为范畴自身的严格对合，即 $\epsilon_{\A}\circ\epsilon_{\A}=\Id$ 逐项成立；而真正等变的结构只应在同伦意义下满足这一条，即带有更高阶的同伦相干数据。两者在上链层次上给出的不动对象并不相同，这正是注 6.5 所说「不协调」的来源。本章只在特征为零或特征 $p\gg0$ 时使用朴素版本，因为那时取不变部分是正合函子（Maschke 定理），不会丢失信息。[来源：原书注 6.5 与 §10b、§14b，并见本章 6(a) 关于 $\operatorname{char}(\K)\neq2$ 的假设。]}

我们使用不动对象的方式也有些朴素。在恰当的范畴表述中，应当改为考虑“等变对象”：它们只在指定的同构意义下被群作用固定。其上同调层次的版本将在后面简要出现（\secref{14b}），但定义适当的上链层次概念并不直接。从形式上说，这种方法的优点在于：可以分别把 $\Tw^{\mathrm{eq}}\A$ 与 $\modcat^{\mathrm{eq}}(\A)$ 定义为 $\Tw\A$ 与 $\modcat(\A)$ 上诱导作用的不动对象子范畴。不过，这里采用的专门定义已经足以满足我们的目的。
\end{note}
